\documentclass[12pt,a4paper]{article}
\usepackage[left=3cm,right=2cm,top=3cm,bottom=2cm]{geometry}
\usepackage[utf8]{inputenc}
\usepackage[T1]{fontenc}
\usepackage[portuguese]{babel}
\usepackage{setspace}
\usepackage{amsmath, amssymb, amsthm, amsfonts}
\usepackage{graphicx}
\usepackage{multirow}
\usepackage[shortlabels]{enumitem}
\usepackage{color}
\usepackage{float}
\usepackage{mathrsfs}
\usepackage{indentfirst}
\usepackage{nopageno}
\usepackage{graphicx}
\usepackage{float}

\setlength{\parindent}{1.25cm}
\renewcommand{\arraystretch}{1.3} % Aumenta a altura das linhas em 1.3 vezes

% === Título === %
\title{Avaliação}
\date{}
\author{\textbf{}}


\begin{document}

\thispagestyle{empty}


% ======= Cabeçalho ======= %
\begin{table}[H]
    \begin{tabular}{c l}
    \multirow{4}{*}{
    \includegraphics[scale=0.13]{img/brasao-uft.png}}
        &  \textbf{UNIVERSIDADE FEDERAL DO TOCANTINS}\\
	    & \textbf{CÂMPUS DE PALMAS}\\
	    & \textbf{CURSO DE MATEMÁTICA}\\ 
	    & \textbf{Disciplina:} Cálculo Diferencial\\
	    & \textbf{Nome:} \rule{0.7\linewidth}{0.2mm}\\
    \end{tabular}
\end{table}
% === Fim do Cabeçalho ===%

\onehalfspacing
\vspace{-2cm}
\textbf{{\let\newpage\relax\maketitle}}
\vspace{-2.5cm}

\begin{enumerate}[leftmargin=*, label=\textbf{\arabic*.}]

    \item Calcule a derivada da seguinte função:
    \[
    f(x) = 3x^4 - 5x^3 + 2x^2 - 7x + 4
    \]
    \vspace{1cm}

    \item Determine os pontos críticos da função e classifique-os como máximos, mínimos ou pontos de sela:
    \[
    g(x) = x^3 - 6x^2 + 9x + 15
    \]
    \vspace{1cm}

    \item Encontre a derivada implícita da seguinte equação:
    \[
    x^2 + y^2 = 25
    \]
    \vspace{1cm}

    \item Qual das alternativas a seguir é a derivada da função \( h(x) = e^{2x} \)?
    \begin{enumerate}[leftmargin=*, label=\textbf{\alph*)}]
        \item \( 2e^{2x} \)
        \item \( e^{2x} \)
        \item \( 2xe^{2x} \)
        \item \( 4e^{2x} \)
    \end{enumerate}
    
    \vspace{1cm}

    \item Temos um pedaço de papelão de 50 cm por 20 cm e vamos recortar os cantos e dobrar as laterais para formar uma caixa. Determine a altura da caixa que dará volume máximo.

    \begin{figure}[H]
        \centering
        \includegraphics[width=0.6\textwidth]{img/caixa_max_volume.png}
        \label{fig:caixa_max_volume}
    \end{figure}

\end{enumerate}

\end{document}


